\begin{kbvrule}
\noindent{\kbvsecfont ಸೂತ್ರ :}\hspace{0.5em}ಸಂಧಿ ಕಾರ್ಯದಲ್ಲಿ ಇ, ಈ, ಉ, ಊ, ಋ-ಕಾರಗಳು ಸವರ್ಣವಲ್ಲದ ಸ್ವರ ಪರವಾದರೆ ಇ, ಈ ಕಾರಗಳಿಗೆ ಯ ಕಾರವೂ ಉ, ಊ ಕಾರಗಳಿಗೆ ವ್‌-ಕಾರವೂ ಋ-ಕಾರಕ್ಕೆ ರ್‌-ಕಾರವೂ ಆದೇಶಗಳಾಗಿ ಬರುವು ದಾದರೆ ಅದು ಯಣ್‌ ಸಂಧಿ ಎನಿಸುವುದು.
\end{kbvrule}

\begin{kbvderiv}
\vspace{4pt plus 3pt minus 1pt}
\begin{center}\begin{adjustbox}{max width=\textwidth}
{\renewcommand{\arraystretch}{1.22}\setlength{\tabcolsep}{4pt}
\begin{tabular}{l c l}
ಗತಿ + ಅಂತರ & {\color{kbvgrey}$=$} & \kbvresult{ಗತ್ಯಂತರ} \\
ಅತಿ + ಅಂತ & {\color{kbvgrey}$=$} & \kbvresult{ಅತ್ಯಂತ} \\
ಆದಿ + ಅಂತ & {\color{kbvgrey}$=$} & \kbvresult{ಆದ್ಯಂತ} \\
ಮನು + ಅಂತರ & {\color{kbvgrey}$=$} & \kbvresult{ಮನ್ವಂತರ} \\
\end{tabular}}
\end{adjustbox}\end{center}
\vspace{4pt plus 3pt minus 1pt}
\end{kbvderiv}

\indent : — ವ್ಯಂಜನ ಸಂಧಿಗಳು\par

\vspace{4pt plus 4pt minus 1pt}\noindent{\kbvsecfont\fontsize{11pt}{14.5pt}\selectfont\color{kbvaccentdark} (೧)\hspace{0.45em}ಜಶ್ತ್ವ ಸಂಧಿ:}\hspace{0.5em}ಕೆಳಗೆ ಕೊಟ್ಟಿರುವ ಉದಾಹರಣೆಗಳನ್ನು ನೋಡೋಣ.\par

\indent ದಿಕ್‌ + ದರ್ಶನ = ದಿಗ್‌ + ದರ್ಶನ = ದಿಗ್ದರ್ಶನ (ಕಕಾರಕ್ಕೆ ಗಕಾರಾದೇಶ)\par

\begin{kbvderiv}
\vspace{4pt plus 3pt minus 1pt}
\begin{center}\begin{adjustbox}{max width=\textwidth}
{\renewcommand{\arraystretch}{1.22}\setlength{\tabcolsep}{4pt}
\begin{tabular}{l c l}
ಅಚ್‌ + ಅಂತ & {\color{kbvgrey}$=$} & ಅಜ್‌ + ಅಂತ = ಅಜಂತ \\
ಷಟ್‌ + ಆನನ & {\color{kbvgrey}$=$} & ಷಡ್‌ + ಆನನ = ಷಡಾನನ \\
\end{tabular}}
\end{adjustbox}\end{center}
\vspace{4pt plus 3pt minus 1pt}
\end{kbvderiv}

\indent (ಟಕಾರಕ್ಕೆ ದಕಾರಾದೇಶ)\par

\indent ಜಗತ್‌ + ಈಶ್ವರ = ಜಗದ್‌ + ಈಶ್ವರ = ಜಗದೀಶ್ವರ (ತಕಾರಕ್ಕೆ ದಕಾರಾದೇಶ)\par

\indent ಅಪ್‌ + ಜ = ಅಬ್‌ + ಜ = ಅಬ್ಜ (ಪಕಾರಕ್ಕೆ ಬಕಾರಾದೇಶ)\par

\indent ಇಲ್ಲಿ ಐದು ಉದಾಹರಣೆಗಳಿವೆ. ಈ ಉದಾಹರಣೆಗಳಲ್ಲಿ ಪೂರ್ವ ಶಬ್ದದ ಕೊನೆಯಲ್ಲಿರುವ ಕ್‌, ಚ್‌, ಟ್‌, ತ್‌, ಪ್‌ ಗಳಿಗೆ ಕ್ರಮವಾಗಿ ಗ್‌, ಜ್‌, ಡ್‌, ದ್‌, ಬ್‌ ಗಳು ಆದೇಶಗಳಾಗಿ ಬಂದು ಸಂಧಿಕಾರ್ಯ ವೇರ್ಪಟ್ಟಿದೆ. ಈ ತೆರನಾದ ಸಂಧಿಯು ಜಶ್ತ್ವ ಸಂಧಿ ಎನಿಸುವುದು.\par
